% Transcription — fonds Alexandre Grothendieck, shelfmark 152, batch 1 (pages 1-12).
% Established from the facsimile archives/batches/152/batch-01.pdf (a mirror of
% https://grothendieck.umontpellier.fr/152.pdf), per the skill
% transcribe-grothendieck. The folder is twelve pages and this is its only batch.
%
% Nine of the twelve leaves are transcribed. Pages 1, 2 and 3 are trial sheets —
% 1 and 2 written landscape in the free bands of a computer listing, 3 a fast
% pencil field on lined paper — and they carry the notion the inventory's title
% names, the « types de mots ». Pages 4, 6, 7, 9, 10 and 12 are the run proper,
% in ink, under his own pagination at the head of the sheet — 1, 1', 2, 2', 3, 3'
% — and under a heading of his own on page 4, « Arbres plans », so the run is
% continuous and complete. Pages 5, 8 and 11 are listings of the same run (IBM
% JCL, the job named PREDONTA) with nothing of his on them, and are skipped; the
% gaps in the page numbers are the record.
%
% The subject is one and the run takes it three times over. Page 2 states the
% equivalence the folder exists for: the category of « types de mots » of length
% n is the isotypic category of systems made of an oriented disc, a « découpage
% par cordes » of it, and at least one marked vertex on each arc that the chords'
% endpoints cut out of the boundary; the structure is written (D, omega, K, S,
% S_0) and its « nombre réduit » is the sum, over the connected components alpha
% of the complement of the chords, of card(alpha inter S) - 1, a type of word
% being trivial when that number is zero and reduced when K is empty. Page 3
% defines admissibility by recursion on the number of edges — sigma = id, or
% there is an orbit of sigma made of two consecutive edges whose contraction
% leaves an admissible graph — and records that every finite graph of order 0 is
% admissible. Page 4 opens the run proper: a tree is a connected 1-connected
% graph, a plane tree is a tree with a class of embeddings, and the embeddings
% are shown to correspond to the circularisations of Gamma, that is to systems of
% circular orders on the sets A_s of edges incident to each vertex. Page 6 cuts
% the surface along |Gamma| and gets a polygon; for a general circularised graph
% the same cutting gives a contour, and the set of arcs Abar carries the
% canonical involution sigma_Gamma reversing orientation, with S isomorphic to
% Abar modulo rho_Gamma. Page 7 names the bijection rd from Abar to the edges of
% C — an arc goes to its right-hand edge, the left-hand one being struck out in
% favour of it — splits Abar(C) into rd(Abar) and rg(Abar), and computes
% sigma_1^C and rho_C on both halves. Page 9 turns that into a dictionary between
% (Abar(Gamma), rho_Gamma, sigma_Gamma) and (A(C), rho_C, sigma_C), and closes on
% a boxed criterion for Gamma to be a tree. Page 10 runs the construction
% backwards — |Gamma| is |C| modulo the gluing of a to sigma_C(a) with
% orientations reversed — and reads the tree case off a disc cut by chords. Page
% 12 states the resulting equivalence of categories.
%
% The hand is drafting throughout and that is the folder's governing fact. Pages
% 4, 9, 10 and 12 read almost end to end; page 6's lower third is a block
% cancelled by two long horizontals and three diagonals over an already crowded
% sheet, and page 7 is a fast palimpsest carrying two marginal notes written in
% biais at some 20 to 25 degrees. Following folder 29 batches 8 and 9, folder 52
% and folder 88, what the leaves will not support is left \ill{} rather than
% invented, and the apparatus is dense in consequence.
%
% Notation fixed once in the header rather than page by page: Abar is his
% notation for the set of arcs, that is of oriented edges, and it is used from
% page 6 on; « rd » and « rg » are his abbreviations for the right-hand and
% left-hand edge of an arc, and the displayed line of page 7 spells the first of
% them out; « circ(Gamma) » is the set of circularisations of Gamma; and « cat. »
% is « catégorie », « long. » « longueur », « nb. » « nombre ». Readings flagged
% and not settled: the noun after « i.e. les \struck{applications} » on page 4,
% read « paires » on the sense and not on the ink, the strokes carrying a tall
% letter the word does not have; the adjective after « graphe » opening the
% Structure of page 3, read « segmenté »; the words closing the boxed criterion
% of page 9; and the subscript of A on the display of page 7, which is drowned in
% ink. One variance of substance is flagged where it occurs and not repaired: the
% head of page 9 writes sigma_Gamma = sigma_C and rho_Gamma = rho_C, and the
% dictionary at the foot of the same page distinguishes rho_Gamma sigma_Gamma
% from rho_Gamma and does not agree with it.
%
% Pass: Opus 5 (claude-opus-5), 2026-09-13 — first pass, unchecked against the pages by a human.
\documentclass[11pt,a4paper]{article}
% Le préambule commun à toutes les transcriptions du fonds Grothendieck.
%
% Il tient en une idée : **l'incertitude se compose, elle ne se résout pas.**
% Une transcription qui lisse un mot douteux a détruit la seule chose pour
% laquelle on la faisait — un lecteur doit pouvoir savoir, à chaque endroit,
% ce qui était sur le feuillet et ce qui a été ajouté. Les six macros qui
% suivent sont donc l'appareil critique, et non de la décoration.
%
% Les mêmes commandes sont reconnues par `scripts/render.mjs`, qui produit la
% vue de lecture du site. Une macro ajoutée ici doit l'être aussi là-bas,
% faute de quoi le rendu échouera bruyamment — ce qui est voulu : un rendu
% silencieusement faux est pire qu'un rendu absent.

\ProvidesPackage{grothendieck}[2026/08/08 Transcriptions du fonds Grothendieck]

\usepackage{amsmath,amssymb,amsthm}
\usepackage{mathrsfs}
\usepackage{stmaryrd}
% Les diagrammes commutatifs sont partout dans ce fonds : c'est la langue
% même de la géométrie algébrique de Grothendieck, et une transcription qui
% les rendrait en prose ne transcrirait rien.
\usepackage{tikz-cd}
% Français par défaut — c'est la langue des feuillets — mais l'anglais est
% chargé aussi : la traduction et le résumé sont des documents anglais, et la
% typographie française (espaces avant les deux-points) y serait une faute.
% Ces documents basculent par \selectlanguage{english} en tête de corps.
\usepackage[english,french]{babel}
\usepackage{fontspec}
\usepackage{geometry}
\geometry{a4paper,margin=25mm,marginparwidth=28mm}
\usepackage{marginnote}
\usepackage{xcolor}
% ulem, not soul. `soul`'s \st and \ul are built on a character-by-character
% scan that XeLaTeX's font machinery breaks: the document fails with an
% undefined control sequence rather than degrading. ulem does the same two jobs
% and is Unicode-engine safe. `normalem` stops it hijacking \emph, which the
% transcriptions use for Grothendieck's own emphasis.
\usepackage[normalem]{ulem}
\usepackage{hyperref}
\hypersetup{colorlinks=true,linkcolor=black,urlcolor=black,pdfborder={0 0 0}}

% --- Métadonnées du lot -----------------------------------------------------
% Reprises telles quelles par le script de rendu, qui les affiche en tête de
% la vue de lecture. Elles fixent ce que le fichier prétend couvrir : sans
% elles, une transcription flotte sans point d'attache dans un fonds de seize
% mille feuillets.

\newcommand{\folder}[1]{\gdef\grothFolder{#1}}
\newcommand{\batch}[1]{\gdef\grothBatch{#1}}
\newcommand{\leaves}[2]{\gdef\grothFirst{#1}\gdef\grothLast{#2}}
\newcommand{\foldertitle}[1]{\gdef\grothFoldertitle{#1}}
\gdef\grothFolder{?}\gdef\grothBatch{?}\gdef\grothFirst{?}\gdef\grothLast{?}\gdef\grothFoldertitle{}

% --- Le résumé --------------------------------------------------------------
%
% Il ouvre la lecture modernisée, sous le titre « Résumé », et il est composé
% en petit. Ce n'est pas un ornement : c'est le seul passage du document qui
% ne soit pas des mathématiques, et un lecteur doit pouvoir le reconnaître
% avant de l'avoir lu — puis le sauter s'il connaît déjà le sujet, ou s'y
% arrêter s'il ne le connaît pas. Le filet qui le suit marque la reprise du
% corps.

\newenvironment{resume}
  {\small\setlength{\parskip}{0.45em}}
  {\par\medskip\hrule\medskip}

% --- Mots-clés ---------------------------------------------------------------
%
% En anglais, en clôture du résumé de la lecture modernisée : le vocabulaire
% moderne sous lequel chercher ce que les feuillets construisent. Le manifeste
% du site (`npm run manifest`) les extrait pour étiqueter la cote entière —
% cette ligne est la source unique des tags, il n'y a pas de fichier à côté.
\newcommand{\keywords}[1]{\par\smallskip\noindent
  {\footnotesize\textsc{Keywords} --- \textit{#1}}\par}

% --- L'appareil critique ----------------------------------------------------

% Le numéro de feuillet, en marge. C'est l'ancre qui permet de régler un
% désaccord sur une formule en regardant *un* feuillet plutôt qu'une chemise
% de six cents. Le site s'en sert aussi pour tourner les pages du fac-similé
% quand on fait défiler la transcription.
\newcommand{\leaf}[1]{%
  \marginnote{\footnotesize\color{gray!70}\textsf{#1}}%
  \label{leaf:#1}\ignorespaces}

% Illisible : on ne devine pas. Un mot inventé qui se lit comme les autres est
% le pire résultat possible.
\newcommand{\ill}{\textcolor{red!60!black}{[\,\ldots\,]}}

% Lecture douteuse : le mot est donné, et signalé comme incertain.
\newcommand{\uncertain}[1]{\uline{#1}}

% Ajout de l'éditeur — une lettre manquante, un mot sous-entendu.
\newcommand{\add}[1]{\textcolor{blue!50!black}{[#1]}}

% Biffé par Grothendieck lui-même. Ce qu'il a rayé fait partie du document :
% une rature dit souvent où la pensée a changé de direction.
\newcommand{\struck}[1]{\sout{#1}}

% Note du transcripteur, sur l'état matériel du feuillet.
\newcommand{\note}[1]{\par\smallskip{\small\color{gray!60!black}\textit{[#1]}}\par\smallskip}

% Note marginale de Grothendieck, distincte du corps du texte.
\newcommand{\marginal}[1]{\par\smallskip\noindent\hspace{1em}%
  {\small\color{gray!60!black}#1}\par\smallskip}

% --- Titre ------------------------------------------------------------------

\AtBeginDocument{%
  \begin{center}
    {\large\bfseries Fonds Alexandre Grothendieck --- cote n\textsuperscript{o}~\grothFolder}\\[2pt]
    {\itshape\grothFoldertitle}\\[4pt]
    {\small feuillets \grothFirst--\grothLast\ (lot \grothBatch)}\\[2pt]
    {\footnotesize\color{gray!60!black}Transcription établie d'après le fac-similé numérisé
     par l'Université de Montpellier.\\
     Les lectures incertaines sont soulignées, les illisibles notées [\,\ldots\,],
     les ajouts entre crochets.}
  \end{center}
  \vspace{1em}}

\endinput


\folder{152}
\batch{1}
\pages{1}{12}
\dating{[à partir de 1982]}
\watermark{Édition de démonstration}
\foldertitle{Catégorie des types de mots [dont arbres plans] : notes manuscrites (s.d.)}

\begin{document}

\section{« Types de mots »}

\note{les trois premiers feuillets sont des feuilles d'essais et n'ont pas de
titre ; les mots mis ici en tête sont les siens, écrits entre guillemets au
milieu de la page 2, et ce sont eux que reprend le titre de l'inventaire. Les
pages 1 et 2 sont écrites en paysage dans les bandes laissées libres par un
listing ; la page 3 est au crayon}

\page{1}

\note{le feuillet est un champ d'essais : une trentaine de petites figures —
des cercles portant des points marqués et des rayons, des arbres à trois ou
quatre branches, deux cercles accolés — parmi lesquelles quelques lignes
seulement sont écrites. Elles sont données ici dans l'ordre où elles se lisent
sur la feuille ; les figures ne sont pas décrites une à une}

\ill{} \quad $S_1$

systèmes \ill{} deux lettres \quad $uv,\ u^{-1}v,\ uv^{-1},\ u^{-1}v^{-1}$

système de trois lettres \quad $uvw$ ou \ill{}

syst.\ $u^{-1}u$ ou $uu^{-1}$

sommets principaux
\note{le mot désigne, sur la figure qui le porte, deux des points marqués sur
le cercle}

$\sigma u^{-1}u$ \qquad $uvu^{-1}$

\note{deux blocs encadrés au bas du feuillet, le premier portant en tête un
mot biffé}

\struck{\ill{}} \quad conjugués \ill{}

\uncertain{invariants} de \ill{}

\note{chacun des deux cadres est suivi d'un petit cercle marqué, le second
partagé par un diamètre}

\ill{} \quad $u+v$ \quad $u+w$

\note{deux essais sont numérotés, 1 et 2, en marge du feuillet ; ils ne sont
suivis de rien}

\page{2}

\marginal{N.B. \ill{} Dans \ill{} \ill{} quelq.\ \ill{} \ill{}
\uncertain{ordonnés}}

La cat.\ des « \struck{\ill{}} types de mots » de long.\ $n$ équivalent :
la cat.\ isotypique des systèmes formés d'un disque \add{orienté} $D$,
d'un « découpage par cordes » de ce disque, avec au moins un sommet sur chaque
arête du \uncertain{contour} découpé sur $\partial D$ par les extrémités des
cordes, \ill{}

\note{« équivalent » est écrit sans « est » ; la phrase reste telle quelle}

$(D,\ \omega,$ \struck{$K$} \struck{$S$} $,\ K,\ S_0)$

\note{les deux lettres biffées sont réécrites aussitôt après, dans l'autre
ordre}

\[
\text{nb.\ d'ordre réduit} \;=\; \sum_{\alpha \in \pi_0(\mathring{D} \smallsetminus K \cap \mathring{D})} \bigl(\operatorname{card}(\alpha \cap S) - 1\bigr)
\]

(ou $\sum_{\alpha \subset P_0}$ \ill{})

\note{le second membre de cette variante est écrit sous l'indice du premier et
n'est pas lisible}

type de mot \uncertain{trivial} : ordre réduit $= 0$

type de mot réduit : $K = \emptyset$ \quad (i.e.\ trivial)

$\operatorname{card} S = n \ (\neq 0)$ \qquad \struck{\ill{}} \quad $\forall\,\alpha$

\note{en haut du feuillet, un tableau de correspondances en trois colonnes,
écrit à la file et sans phrases suivies}

disques \add{ou} découpés par cordes $\longleftrightarrow$ \add{arbres} graphes
circularisés $\longleftrightarrow$ contours \quad $\simeq$ involutions

disques \struck{ou polygones} \struck{découpés}, $\longrightarrow$ graphes
(à bouts libres) circularisés $\longleftrightarrow$ \ill{} contours

une condition que disques \ill{} découpés \ill{}

correspondance \ill{} \uncertain{seulement}, ou une plus \ill{} \ill{}
(aux choix)

\page{3}

\note{le feuillet est occupé aux trois quarts par de grandes figures au
crayon : un disque dont le bord porte des points marqués et qui est découpé par
des cordes, avec l'arbre dual tracé à l'intérieur ; puis deux disques reliés par
une arête, le second plus petit et rayonnant. La prose est écrite au bas et au
milieu, entre les figures}

Structure : graphe \uncertain{segmenté}, une involution $\sigma$ \ill{} sur
l'ens.\ des arêtes.

\uncertain{admissible} se définit par récurrence sur le \ill{} \ill{} d'arêtes
(ordre du graphe)

\begin{quote}
$\Gamma$ admissible d'ordre $n \geq 1$ \quad $\big|$ \quad
$\sigma = \mathrm{id}$

\qquad ou $\exists$ orbite $\varepsilon$ de $\sigma$ formée de deux arêtes
consécutives,

\qquad et le graphe $\Gamma'$ déduit de $\Gamma$ par contraction en $1$ pt
\ill{} \struck{deux} \ill{} des arêtes $(\neq)$, et muni de l'involution
$\sigma'$ déduite de $\sigma$, est admissible
\end{quote}

\note{l'exposant $\varepsilon$ porté par « orbite » est écrit petit et
haut ; il n'est pas repris ailleurs}

Tout graphe fini d'ordre $0$ est admissible \quad (d'ordre $1 \to n$)

\section{Arbres plans}

\note{le titre est le sien, souligné en tête de la page 4, qui porte aussi le
premier chiffre de sa pagination propre. Celle-ci court sur les six feuillets
du texte suivi — $1$, $1'$, $2$, $2'$, $3$, $3'$ sur les pages 4, 6, 7, 9, 10
et 12 des archivistes}

\page{4}\note{sa page 1}

Rappelons qu'un \textbf{arbre} (est un graphe \add{connexe} $1$-connexe (ce qui
implique que c'est un graphe strict, i.e.\ sans boucle fermée, et par deux
\struck{\uncertain{points}} \add{sommets} distincts passe au plus une arête).

\note{« (connexe) » est écrit une seconde fois en interligne au-dessus de la
ligne ; et un « s » est ajouté au-dessus de « arête »}

On peut donc considérer la structure donnée par $(S,A)$, $S = $ ens.\ des
sommets, $A \subset \mathfrak{P}_2(S)$. \struck{\ill{}} L'on a $S \neq \emptyset$,
et $A$ \struck{$\neq$} $= \emptyset \Longleftrightarrow \operatorname{card} S = 1$
(arbre \uncertain{ponctuel}).

On appelle \textbf{arbre plan} un arbre muni d'une classe de plongements
\add{\ill{}} (dans une surface \ill{}) (deux plongements étant identifiés
$|\Gamma| \hookrightarrow X$ \add{et} $|\Gamma| \hookrightarrow X'$ ; il existe
des voisinages \add{$U$, $U'$} de $|\Gamma|$ dans $X$, $X'$ et un homéo.\
$U \simeq U'$ induisant l'identité sur $|\Gamma|$).

Les plongements \ill{} (\ill{} que $|\Gamma|$ est $1$-connexe \add{et} admet un
voisinage \ill{} $1$-connexe $U$ dans $X$, — d'où $U \simeq \mathbb{R}^2$ —)
s'identifient \struck{à} \ill{} aux \textbf{bicircularisations} de $\Gamma$,
i.e.\ les \struck{applications} \uncertain{paires} $(\varepsilon, c)$,
$\varepsilon$ un ens.\ à $2$ éléments \struck{\ill{}} et
$c : \varepsilon \longrightarrow \operatorname{circ}(\Gamma)$ compatible avec le
« \uncertain{passage} à l'opposé ».

\note{le mot lu « paires » porte une haste que le mot n'a pas ; la lecture
suit le sens de la phrase, où $(\varepsilon, c)$ est bien un couple, et non le
tracé. La seconde lettre du couple est écrite par-dessus une autre et vaut le
$c$ de la ligne suivante}

Ici $\operatorname{circ}(\Gamma)$ est l'ens.\ des \textbf{circularisations} de
$\Gamma$, i.e.\ des systèmes d'ordres circulaires sur les $A_s$ (i.e.\ pour
$\forall\, s \in S$, $A_s$ est l'ens.\ des arêtes incidentes à $s$).

Une \struck{plongement} \add{classe de} plongement de $|\Gamma|$ dans une
surface orientée correspond à une circularisation de $\Gamma$ ; \ill{} \ill{}
de l'une \ill{} des $A_s$ $(s \in S)$ telles que $\operatorname{card} A_s \geq 3$.
Donc si les sommets de $\Gamma$ sont d'ordre $\leq 1$, il y a une \textbf{seule}
\ill{}

\page{6}\note{sa page $1'$}

\struck{Si $\Gamma$ n'est pas \ill{} ponctuel,} \struck{Passons à la
description} \ill{}

Prenant le découpage de $X$ \add{suivant $|\Gamma|$}, on trouve une surface avec
un bord connexe, qui est munie de façon naturelle d'une structure de polygone
topologique (i.e.\ \ill{} un ens.\ fini \add{de sommets} \ill{} dessus), et le
polygone combinatoire corresp.\ \ill{} ne dépend que de $\Gamma$, à isom.\ can.\
près. \struck{Voici}

Si on partait d'un graphe circularisé plus général, la \struck{même}
construction marcherait, mais on trouverait un \textbf{contour} \ill{} (pas un
contour connexe), dont voici la construction. Soit $\bar A$ l'ens.\ des
\textbf{arcs} ($=$ arêtes orientées) de $\Gamma$, $\sigma_\Gamma$ l'involution
can.\ sur $\bar A$ (qui \uncertain{renverse} l'orientation \struck{$|\Gamma|$}).

\note{à partir d'ici et jusqu'à la fin du feuillet, un bloc annulé par deux
longues horizontales et trois diagonales, écrit sur une page déjà chargée ; ce
qu'il porte encore est donné, le reste est laissé \ill{}}

\struck{\ill{}} \struck{$\bar A$ s'identifie} \ill{} \struck{aux arêtes}
\ill{} (\ill{} : (note générale)) \ill{} \struck{difficile} \ill{}
\struck{il reste à définir} $\sigma_\Gamma$. \ill{} \struck{l'orientation}
$\rho_\Gamma$ (rendant \ill{}) \ill{} \struck{décrire} \ill{}
\struck{induit par} \ill{}

\[
S \simeq \bar A / \rho_\Gamma
\]

(si $\Gamma$ n'a pas de sommet isolé).

\page{7}\note{sa page 2}

On a une \uncertain{identification} bijective :

\[
\bar A \ \xrightarrow{\ \mathrm{rd}\ }\ A(C) \qquad \text{ens.\ des arêtes de } C
\]

$\bar a \longmapsto$ « \uncertain{arête} \struck{gauche} droite de
$\bar a$ »

\note{l'indice porté par le $A$ du second membre est noyé d'encre et n'est pas
lisible ; la suite du feuillet écrit ce même ensemble $A(C)$ sans indice}

et chaque arête de $C$ est orientée de façon naturelle via cette
\struck{\ill{}} ; $C$ est orienté de façon \ill{}. Donc l'ens.\ des \ill{}
$\bar A(C)$ s'identifie : $A(C) \times \{\pm 1\} \simeq \bar A \times \{\pm 1\}$,
et l'identification $\sigma_0^C$ (\ill{}) : l'autre \struck{\ill{}}
s'identifie : \struck{\ill{}} $\sigma_0 \times \tau$, \ill{}

Reste : définir $\sigma_1^C$, \ill{}

\[
\bar A \times \{+1\} = \bar A(C)_+ \quad \bigl[\, = \mathrm{rd}(\bar A)\,\bigr]
\qquad
\bar A \times \{-1\} = \bar A(C)_- \quad \bigl[\, = \mathrm{rg}(\bar A)\,\bigr]
\]

\note{les deux accolades sont écrites sous les membres de droite ; la seconde
est en partie couverte}

\begin{quote}
$\sigma_1^C(\mathrm{rd}(\bar a)) =$ \struck{\ill{}} $\mathrm{rd}(\rho_\Gamma^{-1}(\bar a)) = \sigma_0^C\,\mathrm{rd}(\rho_\Gamma(\bar a))$

$\sigma_1^C(\mathrm{rg}(\bar a)) = \mathrm{rd}(\rho_\Gamma(\bar a))$ \struck{\ill{}}
\end{quote}

\[
\Longrightarrow \quad
\rho_C(\mathrm{rd}\,\bar a) = \mathrm{rd}(\rho_\Gamma(\bar a)), \qquad
\rho_C(\mathrm{rg}\,\bar a) = \mathrm{rg}(\rho_\Gamma(\bar a))
\]

\note{l'exposant de $\rho_\Gamma$ au second membre de la seconde formule est
noyé d'encre ; les deux premières formules sont encadrées et la seconde des
deux est raturée en son milieu}

\ill{} \ill{} des \ill{} \ill{} \ill{} (\ill{} \ill{} autour de \ill{})
\ill{} i.e.\ \ill{}

\marginal{N.B.\ $\bar A(C)$ \ill{} une \uncertain{réunion} \ill{} contours
\ill{}}

\marginal{N.B.\ La définition \ill{} de $C$ \ill{} contours \ill{} dépend
\ill{} \ill{} $C$ est orienté (\ill{})}

\note{les deux N.B.\ sont écrits dans la marge de gauche, en biais à quelque
vingt-cinq degrés, d'une plume plus rapide ; entre eux, un tripode à trois
flèches sortantes}

\page{9}\note{sa page $2'$}

Ces deux derniers \ill{} permettant, pour le contour $C$, de récupérer un graphe
circularisé $\Gamma$, ayant $\bar A$ comme ens.\ des arêtes orientées, avec
opérations $\sigma_\Gamma$, $\rho_\Gamma$ \ill{} définies ainsi :

\[
\sigma_\Gamma = \sigma_C, \qquad \rho_\Gamma = \rho_C
\]

Envisageons la catégorie des graphes combinatoires circularisés $\Gamma$ (dont
les sommets \ill{} \struck{$(\bar A(\Gamma), \rho_\Gamma, \sigma_\Gamma)$}
\ill{}) : elle est équivalente à celle des contours orientés, \ill{} une
involution sans pt fixe $\sigma_C$ \ill{} sur l'une des \uncertain{arêtes},
soit $C$ $(= A(C))$ un \ill{} des arêtes de $S$, $\rho_C$ \ill{} relations
fondamentales des $C$, $\sigma_C$ \ill{} $A(C)$), \struck{\ill{}} $\rho_C$
compatibles \ill{}

\[
\begin{array}{rcl}
\bar A(\Gamma) & \longleftarrow & A(C)\\[2pt]
\sigma_\Gamma & \longleftrightarrow & \rho_C\\[2pt]
\rho_\Gamma \sigma_\Gamma & \longleftarrow & \rho_C\\[2pt]
\rho_\Gamma & \longleftarrow & \rho_C\,\sigma_C
\end{array}
\]

\note{ce dictionnaire ne s'accorde pas avec les deux égalités données en tête
du même feuillet, qui posent $\sigma_\Gamma = \sigma_C$ et
$\rho_\Gamma = \rho_C$ ; les deux états sont donnés tels quels et non
raccordés. La première flèche porte un $\simeq$ au-dessus d'elle}

\marginal{On appelle \ill{} \ill{} \ill{} par \ill{}}

\marginal{$\rho_\Gamma$ relativement \ill{} pour la carte \ill{} définie par
$\Gamma$}

\begin{quote}
$\Gamma$ un arbre $\Longleftrightarrow$ $C$ est un polygone connexe (non
ponctuel), et l'involution $\sigma_C$ dans $A(C)$ est totalement non \ill{},
i.e.\ pour deux \struck{\ill{}} \add{orbites} disjointes d'arêtes sous
$\sigma_C$, \struck{la} composition est \ill{}
\end{quote}

\note{l'énoncé est encadré au bas du feuillet ; les deux derniers mots ne se
laissent pas lire}

\page{10}\note{sa page 3}

\begin{quote}
\add{(non pointés)}

« Arbres plans » \quad (\ill{} \ill{} plus pas circ.\ \struck{orienté})

\qquad $\updownarrow \simeq$

polygones combinatoires \ill{}, une \struck{relation} (ou \ill{}) \ill{}
\ill{} sur l'une des \struck{sommets} arêtes.
\end{quote}

\note{ce bloc est écrit en tête du feuillet et séparé de la suite par un
trait}

Revenons au cas d'un graphe $\Gamma$ circularisé \add{général}. Topologiquement,
le graphe top.\ $|\Gamma|$ se déduit de $|C|$ par
\struck{les identifications \ill{}} \ill{} \ill{} quotient par la relation
d'équivalence suivante : on recolle l'arête $a$ et l'arête $\sigma_C^a$, en
\uncertain{renversant} les \struck{orientations} \add{induites} \ill{}
\uncertain{orientations}.

\note{il écrit l'involution en exposant, $\sigma_C^a$, là où le reste du
dossier l'écrit $\sigma_C(a)$ ; c'est transcrit comme il l'a laissé}

Dans le cas où $\Gamma$ est un arbre, on peut interpréter autrement — c'est
\ill{} typique considéré \ill{} $C$ comme \ill{} d'un disque orienté
\struck{\ill{}} en arêtes, et en joignant les « cordes » \ill{} \ill{} qui se
correspondent par $\sigma_C$, de façon à ce que les \ill{} obtenus ne se
rencontrent pas.

\note{une figure au bas du feuillet : un disque dont le bord porte des points
marqués, et trois cordes tracées à l'intérieur}

\page{12}\note{sa page $3'$}

\ill{} le graphe $\Gamma$ expriment la façon dont les régions se recollent
suivant les frontières communes.

\add{(pas nécess.\ plus orientés)}

La catégorie des \struck{arbres plans} est équivalente : la catégorie
isotypique des disques \struck{\ill{}} (pas nécess.\ orientés de \ill{}) munis
d'un découpage.

Le polygone combinatoire \ill{} de celui \struck{\ill{}} défini par les cordes
du découpage \ill{} sur le bord des disques, les sommets de $\Gamma$
correspondent à des paquets d'arêtes des \ill{} (celles qui bordent une \ill{}
du découpage \ill{}).

\note{le feuillet s'arrête là ; le dossier ne va pas plus loin}

\end{document}
